%% -!TEX root = AUTthesis.tex
% در این پرونده، عنوان پایان‌نامه، مشخصات خود، متن تقدیمی‌، ستایش، سپاس‌گزاری و چکیده پایان‌نامه را به فارسی، وارد کنید.
% توجه داشته باشید که جدول حاوی مشخصات پروژه/پایان‌نامه/رساله و همچنین، مشخصات داخل آن، به طور خودکار، درج می‌شود.
%%%%%%%%%%%%%%%%%%%%%%%%%%%%%%%%%%%%
% دانشکده، آموزشکده و یا پژوهشکده  خود را وارد کنید
\faculty{دانشکده مهندسی کامپیوتر}
% گرایش و گروه آموزشی خود را وارد کنید
\department{مهندسی کامپیوتر}
% عنوان پایان‌نامه را وارد کنید
\fatitle{تشخیص توضیح‌پذیر سلامت روان با بهره‌گیری از مدل‌های زبانی بزرگ و گراف‌های دانش دامنه‌ای}
% نام استاد(ان) راهنما را وارد کنید
\firstsupervisor{دکتر محمد رحمتی}
%\secondsupervisor{استاد راهنمای دوم}
% نام استاد(دان) مشاور را وارد کنید. چنانچه استاد مشاور ندارید، دستور پایین را غیرفعال کنید.
%\secondadvisor{استاد مشاور دوم}
% نام نویسنده را وارد کنید
\name{علیرضا }
% نام خانوادگی نویسنده را وارد کنید
\surname{هوشیاری}
%%%%%%%%%%%%%%%%%%%%%%%%%%%%%%%%%%
\thesisdate{شهریور ۱۴۰۵}

% چکیده پایان‌نامه را وارد کنید
\fa-abstract{
افزایش شیوع اختلالات روانی، ضرورت توسعه ابزارهای دقیق، قابل اعتماد و توضیح‌پذیر برای پشتیبانی از فرایند تشخیص را برجسته ساخته است. مدل‌های زبانی بزرگ می‌توانند هم‌زمان برچسب تشخیصی و توضیح متنی تولید کنند؛ با این حال، توهم اطلاعاتی و فقدان پشتوانه دانشی صریح، استفاده از آن‌ها را در حوزه حساس سلامت روان با چالش مواجه می‌کند. هدف این پژوهش، بازپیاده‌سازی، انطباق و مقایسه روش‌های گراف‌محور برای تشخیص توضیح‌پذیر سلامت روان است. روش‌های منتخب از سه الگوی بازیابی‌محور، عامل‌محور و قیدگذارانه استفاده می‌کنند و نقش گراف دانش دامنه‌ای در پشتیبانی از استدلال و کاهش ادعاهای بدون پشتوانه در آن‌ها بررسی می‌شود. آزمایش‌ها بر روی وظایف منتخب مجموعه‌داده \lr{IMHI} انجام خواهند شد و روش‌ها از نظر درستی پیش‌بینی، کیفیت و اتکاپذیری توضیح، زمان پاسخ و هزینه متنی مقایسه می‌شوند. همچنین، توضیحات تولیدشده با استفاده از چارچوب \lr{GraphEval} در سطح ادعا و روابط ساختاری تحلیل خواهند شد.
}



% کلمات کلیدی پایان‌نامه را وارد کنید
\keywords{
	سلامت روان،
	تشخیص توضیح‌پذیر،
	مدل‌های زبانی بزرگ،
	گراف دانش،
	ارزیابی گراف‌محور
}



%%%%%%%%%%%%%%%%%%%%%%%%%%%%%%%%%%
